\section{Problem Summary}
You are given a binary grid where land cells form islands through four-directional adjacency. Count how many separate islands exist.

The hard part is not detecting land. It is making sure an entire connected component is counted exactly once.

\section{Recognition Pattern}
This is a \textbf{DFS/BFS + visited} problem on a grid.

In interviews, recognize it when:
\begin{itemize}
\item the input is a 2D grid
\item cells connect through fixed directions such as up, down, left, and right
\item the question asks for connected components, reachability, or flood fill
\end{itemize}

\section{Key Idea}
Scan the grid cell by cell. Whenever you find unvisited land, that land must belong to a new island, so increment the answer and flood-fill the entire component.

The flood fill marks every cell in that island as visited. After that, the outer scan can safely ignore those cells, so each island contributes exactly one count.

For this problem, mutating the grid from `'1'` to `'0'` is the simplest visited-marking strategy.

\section{Step-by-step Reasoning}
The naive mistake is to count every land cell, but that clearly overcounts multi-cell islands.

The next observation is structural: two land cells belong to the same island precisely when they are in the same connected component under four-directional movement.

That reframes the task as a standard graph problem. The grid is just an implicit graph where each land cell has up to four neighbors. Once you see that, DFS or BFS is the natural tool: start at one land cell, visit everything reachable from it, and treat that whole region as one island.

\section{Complexity}
\begin{itemize}
\item Time: \(O(mn)\)
\item Space: \(O(mn)\) in the worst case due to the recursion stack
\end{itemize}

\section{Common Pitfalls}
\begin{itemize}
\item Only four-directional neighbors count. Diagonals do not connect islands here.
\item Mark a cell as visited before exploring neighbors, otherwise cycles cause repeated work.
\item Mutating the grid is clean here, but in other interview settings you may need a separate `visited` matrix if the input must stay unchanged.
\end{itemize}

\section{Variants / Follow-ups}
This pattern leads directly to \textbf{Max Area of Island}, \textbf{Surrounded Regions}, \textbf{Flood Fill}, and \textbf{Number of Provinces}.

The reusable idea is to translate a grid into component traversal instead of treating each cell as an isolated local decision.
